\section{Gauge normal form and filtered principal-part obstructions}
\label{sec:filtered-principal-parts}

Retain the notation
\[
 A=A_0,\qquad B=B_0,\qquad
 \Omega=2A\,dB-3B\,dA=z^2\,dz.
\]
For a layer variation $(a,b)$ put
\[
 u=\frac aA,\qquad v=\frac bB,\qquad
 \xi=u+v,\qquad \eta=2v-3u.
\]

\begin{proposition}[Gauge normal form]
\label{prop:gauge-normal-form}
For every normal order $r$,
\begin{equation}
\label{eq:gauge-normal-form}
 \boxed{
 \mathscr D_r(a,b)
 =AB\left(d\eta+\frac r5\eta\,d\log(AB)\right)
 +\frac{5-r}{5}\xi\,\Omega.}
\end{equation}
Consequently, over the rational function field of the boundary curve,
$\mathscr D_r$ is surjective for $r\ne5$.  At the resonant order,
\begin{equation}
\label{eq:resonant-exact}
 \boxed{\mathscr D_5(a,b)=d(2Ab-3Ba).}
\end{equation}
\end{proposition}

\begin{proof}
Substitute
\[
 u=\frac{2\xi-\eta}{5},\qquad
 v=\frac{3\xi+\eta}{5}
\]
into
\[
 \mathscr D_r(a,b)
 =(2-r)a\,dB-3B\,da+2A\,db+(r-3)b\,dA
\]
to obtain \eqref{eq:gauge-normal-form}.  If $r\ne5$, set $\eta=0$ and
$\xi=5\Phi/((5-r)\Omega)$ to solve $\mathscr D_r(a,b)=\Phi$ over the
function field.  For $r=5$, the $\xi$ term disappears and
$AB(d\eta+\eta\,d\log(AB))=d(AB\eta)=d(2Ab-3Ba)$.
\end{proof}

\begin{remark}[An extended deformation complex]
Allowing a variation $\rho\Omega$ of the normalized right-hand side, put
\[
 \widetilde{\mathscr D}_r(a,b,\rho)
 =\mathscr D_r(a,b)-\rho\Omega.
\]
Then
\[
 G_r(h)=(2Ah,3Bh,(5-r)h)
\]
satisfies $\widetilde{\mathscr D}_r\circ G_r=0$.  Thus the fixed layer is
represented by the three-term complex
\[
 \mathcal G_r\xrightarrow{G_r}
 \mathcal U_{A,r}\oplus\mathcal U_{B,r}\oplus\mathcal R_r
 \xrightarrow{\widetilde{\mathscr D}_r}\mathcal W_r.
\]
At $r=5$ the right-hand-side component of $G_r$ vanishes, leaving the
residual normal-coordinate gauge in the fixed-volume problem.
\end{remark}

\subsection{The exact two-point windows}

Let $o=(z=0)$ and $p=(z=\infty)$ on $E\simeq\mathbf P^1$.  For
$5\le r\le8$, the full Newton support gives
\begin{align}
 U_{A,r}&=\langle1,z,\ldots,z^{10-r}\rangle
        =H^0(E,\mathcal O((10-r)p)),\label{eq:UA-window}\\
 U_{B,r}&=\langle1,z,\ldots,z^{15-r}\rangle
        =H^0(E,\mathcal O((15-r)p)),\label{eq:UB-window}\\
 W_r&=\langle z^{-1},1,z,\ldots,z^{24-r}\rangle\,dz
     =H^0(E,\Omega_E^1(o+(26-r)p)).\label{eq:W-window}
\end{align}
The exact layer matrices have full column rank, hence
\[
 \dim\operatorname{coker}\mathscr D_r
 =(26-r)-(11-r)-(16-r)=r-1.
\]

\begin{proposition}[Filtered principal-part duality]
\label{prop:filtered-principal-part-duality}
Order the target basis as $z^I dz$, $-1\le I\le24-r$.  If
$c=(c_I)$ is a row vector, set
\[
 \lambda_c(z)=\sum_{I=-1}^{24-r}c_Iz^{-I-1}.
\]
Then $c$ lies in the left kernel of the layer matrix if and only if
\begin{align}
 \operatorname{Res}_{0}\!
 \left(z^i\bigl(3B\,d\lambda_c+(5-r)\lambda_c\,dB\bigr)\right)&=0,
 &&0\le i\le10-r,\label{eq:filtered-adjoint-A}\\
 \operatorname{Res}_{0}\!
 \left(z^j\bigl(-2A\,d\lambda_c+(r-5)\lambda_c\,dA\bigr)\right)&=0,
 &&0\le j\le15-r.\label{eq:filtered-adjoint-B}
\end{align}
For the lower-layer forcing form $\Phi_r$, the corresponding compatibility
coordinate is
\begin{equation}
\label{eq:residue-coordinate}
 c^T(-\Phi_r)
 =\operatorname{Res}_{0}\bigl(\lambda_c(-\Phi_r)\bigr).
\end{equation}
\end{proposition}

\begin{proof}
The identity
\begin{align*}
 \lambda\mathscr D_r(a,b)
 &-a\bigl(3B\,d\lambda+(5-r)\lambda\,dB\bigr)\\
 &-b\bigl(-2A\,d\lambda+(r-5)\lambda\,dA\bigr)
 =d(-3\lambda Ba+2\lambda Ab).
\end{align*}
shows that the residue of the left side is zero.  Taking successively
$a=z^i,b=0$ and $a=0,b=z^j$ gives
\eqref{eq:filtered-adjoint-A}--\eqref{eq:filtered-adjoint-B}.  Conversely,
those equalities say exactly that the residue functional annihilates every
column of the ordered coefficient matrix.  Formula
\eqref{eq:residue-coordinate} is the definition of the dual monomial basis.
\end{proof}

\begin{remark}
Away from resonance the two adjoint expressions in
\eqref{eq:filtered-adjoint-A}--\eqref{eq:filtered-adjoint-B} need not vanish
as meromorphic differentials.  Only their projections to the duals of the
allowed Newton windows vanish.  The obstruction classes are therefore
filtered principal parts, rather than unrestricted global adjoint
solutions.
\end{remark}

\subsection{The exact cokernel frames through order eight}

The deterministic row reduction gives echelon principal-part frames whose
leading pole orders at $o$ are
\[
\begin{array}{c|c|c}
 r&\dim\operatorname{coker}\mathscr D_r&
 \text{leading pole orders}\\
\hline
5&4&0,18,19,20\\
6&5&0,16,17,18,19\\
7&6&0,14,15,16,17,18\\
8&7&0,12,13,14,15,16,17.
\end{array}
\]
In every layer the order-zero class is represented by $\lambda=1$ and its
forcing pairing vanishes identically.  At $r=5$, the order-twenty class also
pairs identically to zero; the order-nineteen pairing vanishes after the
order-four normalization.  At $r=6$, the order-nineteen pairing vanishes
after that normalization.  The surviving normalized equations are therefore
\[
1,3,5,6
\]
at orders $5,6,7,8$, respectively.

In particular, the five weight-seven equations are exactly the residue
coordinates with leading pole orders $14,15,16,17,18$, and the six
weight-eight equations are exactly those with leading pole orders
$12,13,14,15,16,17$.  The exact principal parts and normalization scalars
are contained in the accompanying machine-readable certificate.

\begin{remark}[The six resultant coordinates]
The sparse resultant and norm certificate uses
\[
 (F_4,F_6,F_8,F_9,F_{10},F_{11}).
\]
In the filtered frames above these are the layer-seven coordinates of
leading pole orders $14,16,18$ and the layer-eight coordinates of leading
pole orders $12,13,14$.  This is an exact coordinate projection of the full
fifteen-component obstruction section.  It is not a whole layer or a
subspace defined by one pole cutoff.  The present calculation therefore
does not identify this six-dimensional projection as a canonical
subquotient of the pole filtration; its current justification is its exact
unit-ideal and mixed-volume effectiveness.
\end{remark}
